\documentclass{article}
\usepackage[utf8]{inputenc}
\usepackage{hyperref}
\usepackage{graphicx}
\usepackage{listings}
\usepackage{xcolor}

\definecolor{codegreen}{rgb}{0,0.6,0}
\definecolor{codegray}{rgb}{0.5,0.5,0.5}
\definecolor{codepurple}{rgb}{0.58,0,0.82}
\definecolor{backcolour}{rgb}{0.95,0.95,0.92}

\lstdefinestyle{mystyle}{
    backgroundcolor=\color{backcolour},   
    commentstyle=\color{codegreen},
    keywordstyle=\color{magenta},
    numberstyle=\tiny\color{codegray},
    stringstyle=\color{codepurple},
    basicstyle=\ttfamily\footnotesize,
    breakatwhitespace=false,         
    breaklines=true,                 
    captionpos=b,                    
    keepspaces=true,                 
    numbers=left,                    
    numbersep=5pt,                  
    showspaces=false,                
    showstringspaces=false,
    showtabs=false,                  
    tabsize=2
}

\lstset{style=mystyle}
\usepackage[a4paper, left = 3cm, right = 3cm, top =2cm]{geometry}
\title{\textbf{\underline{Proyecto de Matemática Numérica}} \\ \textit{Predicción de crecimiento demográfico.} \\ \textit{Informe \#2}}
\author{Luis Ernesto Serras Rimada \\ Guillermo Cepero García \\ Miguel Vadim Vilariño Pedraza}

\date{\today}
\begin{document}
\maketitle
    Informe Plantilla
    \tableofcontents
    \lstlistoflistings
    Introduccion
\newpage
\section{Desarrollo}
Desarrollo\\
\begin{center}
\includegraphics[height = 15cm]{img/cuba3.jpeg}
\end{center}
\newpage
\subsection{Code}
\begin{lstlisting}[language=Python, caption=Python example]
    import numpy as np
        
    def incmatrix(genl1,genl2):
        m = len(genl1)
        n = len(genl2)
        M = None #to become the incidence matrix
        VT = np.zeros((n*m,1), int)  #dummy variable
        
        #compute the bitwise xor matrix
        M1 = bitxormatrix(genl1)
        M2 = np.triu(bitxormatrix(genl2),1) 
    
        for i in range(m-1):
            for j in range(i+1, m):
                [r,c] = np.where(M2 == M1[i,j])
                for k in range(len(r)):
                    VT[(i)*n + r[k]] = 1;
                    VT[(i)*n + c[k]] = 1;
                    VT[(j)*n + r[k]] = 1;
                    VT[(j)*n + c[k]] = 1;
                    
                    if M is None:
                        M = np.copy(VT)
                    else:
                        M = np.concatenate((M, VT), 1)
                    
                    VT = np.zeros((n*m,1), int)
        
        return M
    \end{lstlisting}
\clearpage

\section{Conclusiones}

\cleardoublepage
\section{Referencias}
\url{https://github.com/LFrench03/Modelo-de-Crecimiento-Poblacional/blob/main/data/csv/poblacion-residente.csv}
\\\textit{\textbf{\underline{Referencias}}:}
\begin{itemize}
    \item (1)
\end{itemize}

\end{document}